\sisetup{
    detect-all                  = true,                         % keep formatting around
    separate-uncertainty        = true,                         % 7.2 ± 0.5
    multi-part-units            = brackets,                     % 12.3 ± 0.2 kg, do not repeat "kg"
    per-mode                    = symbol,                       % `symbol` for "m/s", `reciprocal` for "ms^{-1}"
    group-separator             = {\,},
    group-minimum-digits        = 5,
    inter-unit-product          = {},                           % nothing between units
    input-decimal-markers       = {.},
    exponent-product            = \cdot,                        % \times for 5×10^7, \cdot for 5.1O^7
    number-unit-product         = {\ },                         % a non-breakable space between numbers and units
    range-units                 = repeat,
    range-phrase                = {\text{ -- }},                % 5 -- 10 m/s
    list-separator              = {\text{,}\allowbreak\ },
    list-units                  = repeat,                       %
    forbid-literal-units        = true,
}

\DeclareSIUnit{\permille}{\text{\textperthousand}}
\DeclareSIUnit{\ppm}{ppm}
\DeclareSIUnit{\ppb}{ppb}

\DeclareSIUnit{\gforce}{\textit{g}}
\DeclareSIUnit{\galileo}{Gal}                                % gravimetry's cm/s^2
\DeclareSIUnit{\year}{a}

\DeclareSIUnit{\inch}{in}
\DeclareSIUnit[number-unit-product={}]{\inchp}{{}^{\prime\prime}}
\DeclareSIUnit{\foot}{ft}
\DeclareSIUnit[number-unit-product={}]{\footp}{{}^{\prime}}
\DeclareSIUnit{\mile}{mi}

\DeclareSIUnit{\au}{au}
\DeclareSIUnit{\parsec}{pc}
\DeclareSIUnit{\lightyear}{ly}
\DeclareSIUnit{\watthour}{Wh}
\DeclareSIUnit{\jansky}{Jy}
\DeclareSIUnit{\upmag}{\kern-0.25em\ensuremath{^\mathrm{m}}}
\DeclareSIUnit{\mag}{mag}
\DeclareSIUnit{\rpm}{rpm}
\DeclareSIUnit{\textarcsecond}{as}
\DeclareSIUnit{\solarmass}{\ensuremath{\mathrm{M}_{\Sun}}}
\DeclareSIUnit{\solarluminosity}{\ensuremath{\mathrm{L}_{\Sun}}}
\DeclareSIUnit{\earthmass}{\ensuremath{\mathrm{M}_{\Earth}}}
\DeclareSIUnit{\horsepower}{hp}

\DeclareSIUnit{\byte}{B}
\DeclareSIUnit{\litre}{l}
\DeclareSIUnit{\celsius}{\text{°}C}
\DeclareSIUnit{\dcelsius}{\text{°}C}
\DeclareSIUnit{\fahrenheit}{\text{°}F}
% siunitx ships its own `\degreeCelsius`, and it sets a math ring operator with a space before
% the C -- `25 ∘ C` where every other temperature in the booklet is `25 °C`. It is redeclared
% here rather than avoided because pint emits it: up to 0.26.1 a Celsius quantity formatted as
% `\degree_Celsius`, which is not valid TeX, so `PhysicsQuantity.PINT_TO_SIUNITX` caught it by
% name and wrote `\celsius`; 0.26.2 emits `\degreeCelsius`, which is valid TeX and therefore
% never reaches that mapping. 146 source files declare a Celsius quantity.
\DeclareSIUnit{\degreeCelsius}{\text{°}C}
\DeclareSIUnit{\pound}{lb}

\DeclareSIUnit{\calorie}{cal}
% a tonne of TNT equivalent, 4.184 GJ by definition. Prefixable like any other unit, so
% `\mega\tontnt` is the megaton a blast is usually quoted in.
\DeclareSIUnit{\tontnt}{t\text{ TNT}}
\DeclareSIUnit{\dyne}{dyn}
\DeclareSIUnit{\erg}{erg}

\DeclareSIUnit{\pixel}{px}
\DeclareSIUnit{\eur}{\text{€}}
\DeclareSIUnit{\atmosphere}{atm}
\DeclareSIUnit{\torr}{Torr}
\DeclareSIUnit{\atomicmass}{u}

\DeclareSIUnit{\molar}{M}
